% ============================================================
% 分片：s06-categories  对应原文 PDF 第 76--87 页（印刷页 76--87）
% 负责章节：§6 Categories Of Eigensheaves + 6.1--6.6（自 §6 标题起，至 §7 标题前）
% 译者：trans-s06   审校状态：初稿
% ============================================================

% ==== 原文 PDF 第 76 页（印刷页 76）====
% [跨页接续: 原文第 75 页/76 页。本页前两段属 §5.3 结尾，见上一分片 s05b；本分片自 §6 标题起]

\section{本征层的范畴}

\begin{keytakeaway}
\textbf{核心物理图像：}第 5 节在奇异 Hitchin 纤维上得到的 A-膜（特别是分数膜），经 A-膜/D-模对应应给出 $\Bun_G$ 上满足某一推广 Hecke 性质的 D-模；内窥点处 ${}^LG$-局部系统的自同构群 $\Gamma$（有限群）正是控制该 D-模分裂方式的对称性。
\textbf{数学工具：}两类范畴——带群 $H$ 作用的类型 (1) 与带表示张量范畴 $\mathrm{Rep}(H)$ 幺半作用的类型 (2)——以及二者之间的等变化与去等变化；Hecke 函子 $H_V$ 与 Hecke 本征对象；正则 Hecke 性质与分数 Hecke 性质；D-模的 Fourier--Mukai 变换；几何 Eisenstein 级数。
\textbf{主要结论：}在内窥例 $G=SL_2$、${}^LG=SO_3$、$\Gamma=\mathbb{Z}_2$ 中，正则 Hecke 本征层分裂为 $\mathcal{F}=\mathcal{F}_+\oplus\mathcal{F}_-$，两个不可约直和项无典范区分；分数 Hecke 本征层满足 (6.11)，而范畴 $Aut_{\mathcal{E}}$ 虽等价于 $\mathrm{Vect}$，却有两种不可区分的等价。对 $\mathbb{C}^\times$、任意环面、以中心为自同构群以及几何 Eisenstein 级数，给出同一框架下的平行结论。
\end{keytakeaway}

在前几节中，我们在 Higgs 丛的模空间上构造了 A-膜，它们满足一个与通常的 Hecke 性质非常相似、但并不完全相同的性质。如第 4 节所解释的，对每个这样的 A-膜，我们应当能附上一个 $\Bun_G$ 上的 D-模。这些 D-模于是应当满足同样的 Hecke 性质。本节我们解释一个一般框架，在其中可以把这一性质解释为 Hecke 本征层这一标准概念的自然推广。

\subsection{范畴的一般性说明}

设 $H$ 是 $\mathbb{C}$ 上的约化代数群，$\mathrm{Rep}(H)$ 是 $H$ 的有限维表示构成的张量范畴。下面我们考虑与 $H$ 相关联的两类阿贝尔范畴：

（1）带有 $H$ 的作用的范畴。这意味着每个 $h\in H$ 给出范畴 $\mathcal{C}$ 上的一个函子 $F_h$，它把 $\mathcal{C}$ 的每个对象 $M$ 送到对象 $F_h(M)$，把每个态射 $f : M\to M'$ 送到态射 $F_h(f) : F_h(M)\to F_h(M')$。

这些函子必须满足自然的相容性；特别地，对任意 $h,h'\in H$，存在同构 $i_{h,h'} : F_{hh'}\simeq F_h\circ F_{h'}$，并且有 $i_{hh',h''}i_{h,h'} = i_{h,h'h''}i_{h',h''}$。

% [Tutor Note: 推导补足。这条余圈恒等式两边都是 $F_{hh'h''}$ 到 $F_h\circ F_{h'}\circ F_{h''}$ 的同构：左端先用 $i_{hh',h''}$ 把 $F_{hh'h''}$ 送到 $F_{hh'}\circ F_{h''}$，再把 $i_{h,h'}$ 与 $F_{h''}$ 作张量积；右端先用 $i_{h,h'h''}$ 送到 $F_h\circ F_{h'h''}$，再把 $F_h$ 与 $i_{h',h''}$ 作张量积。原文省略了这步 whiskering，所以字面上并不是复合。它正是函子 $h\mapsto F_h$ 的结合约束；缺了它，$h\mapsto F_h$ 只是一个函子族而非作用，$\mathcal{C}$ 也就不成其为类型 (1) 的范畴。]

例：带有 $H$ 的作用的代数簇 $X$ 上凝聚层（或 B-膜）的范畴 $\mathrm{Coh}(X)$。

（2）带有张量范畴 $\mathrm{Rep}(H)$ 的幺半作用的范畴。这意味着任何对象 $V\in\mathrm{Rep}(H)$ 定义 $\mathcal{C}$ 上的一个函子 $M\mapsto V\star M$，并且这些函子与 $\mathrm{Rep}(H)$ 上的张量结构相容。

例：商 $X/H$ 上凝聚层的范畴 $\mathrm{Coh}(X/H)$，其中 $X$ 如上面第（1）点。\footnote{这里为简单起见可以假设该作用是自由的；但若我们愿意把 $X/H$ 视为代数叠——或在 $H$ 为有限群时视为轨形——则这一假设并非必要。}于是对每个 $V\in\mathrm{Rep}(H)$，我们有 $X/H$ 上与主 $H$-丛 $X\to X/H$ 相联系的向量丛 $\mathcal{V}$，
\[
\mathcal{V} = X\underset{H}{\times} V.
\]
其截面层（也记作 $\mathcal{V}$）是 $\mathrm{Coh}(X/H)$ 的对象。对 $\mathrm{Coh}(X/H)$ 的任何其他对象 $M$，我们置
\[
V\star M := \mathcal{V}\underset{\mathcal{O}_{X/H}}{\otimes} M.
\]

可以在这两类范畴之间互相转换：

由（1）到（2）的转换称为等变化（equivariantization）；

由（2）到（1）的转换称为去等变化（de-equivariantization）。

这两个程序互逆。这里我们只作简要勾勒，因为这部分材料相当为人熟知（例如见 [AG]）。

% ==== 原文 PDF 第 77 页（印刷页 77）====

等变化如下进行：给定类型（1）的范畴 $\mathcal{C}$，我们构造类型（2）的新范畴 $\mathcal{C}'$。其对象是数据 $(M,(\mu_h)_{h\in H})$，其中 $M\in\mathcal{C}$ 且
\[
\mu_h : F_h(M)\xrightarrow{\sim} M,\qquad h\in H,
\]
是一族同构，满足 $\mu_{hh'} = \mu_h\circ\mu_{h'}$。给定 $V\in\mathrm{Rep}(H)$，我们定义 $\mathcal{C}'$ 的一个新对象
\[
V\star(M,(\mu_h)_{h\in H}) = (\underline{V}\otimes_{\mathbb{C}} M,(\rho(h)\otimes\mu_h)_{h\in H}),
\]
其中 $\underline{V}$ 是 $V$ 底下的向量空间，$\rho : H\to\mathrm{End}(\underline{V})$ 是 $H$ 在 $V$ 上的表示。于是我们看到，$\mathrm{Rep}(H)$ 作用在 $\mathcal{C}'$ 上。

例：若 $\mathcal{C} = \mathrm{Coh}(X)$，则 $\mathcal{C}'$ 是 $X$ 上 $H$-等变凝聚层的范畴，它与 $X/H$ 上凝聚层的范畴 $\mathrm{Coh}(X/H)$ 相同。

去等变化是类似的，而「Hecke 本征对象」这一概念正是在这里自然地出现。

给定类型（2）的范畴 $\mathcal{C}'$，我们定义类型（1）的新范畴 $\mathcal{C}$。其对象是数据 $(M,(\alpha_V)_{V\in\mathrm{Rep}(H)})$，其中 $M\in\mathcal{C}'$ 且
\begin{equation}\tag{6.1}
\alpha_V : V\star M\xrightarrow{\sim}\underline{V}\otimes_{\mathbb{C}} M,\qquad V\in\mathrm{Rep}(H),
\end{equation}
是一族与 $\mathrm{Rep}(H)$ 上的张量积结构相容的同构。我们把 $(M,(\alpha_V)_{V\in\mathrm{Rep}(H)})$ 称为范畴 $\mathcal{C}'$ 的一个 Hecke 本征对象（Hecke eigenobject）。

% [Tutor Note: 易错预警。(6.1) 右端的 $\underline{V}$ 是表示 $V$ 底下的纯向量空间，$\otimes_{\mathbb{C}}$ 用的是 $\mathcal{C}$ 的 $\mathbb{C}$-线性结构，而不是 $\star$ 作用。因此 (6.1) 并非断言 $V\star M$ 与 $V\otimes M$ 同构；它断言的是：$\star$ 作用在 $M$ 上等于按表示矩阵 $\rho$ 作用于标量因子后所得的张量积。这正是「本征」二字的含义，也是 (6.2) 能用 $\rho(h)$ 扭转同构的原因。]

群 $H$ 自然地作用在 $\mathcal{C}$ 上：对 $h\in H$ 我们定义
\begin{equation}\tag{6.2}
F_h((M,(\alpha_V)_{V\in\mathrm{Rep}(H)})) = (M,((\rho(h)\otimes\mathrm{id})\circ\alpha_V)_{V\in\mathrm{Rep}(H)}).
\end{equation}
换言之，$M$ 保持不变，但我们用 $h$ 扭转同构 $\alpha_V$。于是 $\mathcal{C}$ 确实是类型（1）的范畴。

\subsection{例子}

类型（2）范畴的最简单例子是带有在自身上的自然幺半作用的范畴 $\mathrm{Rep}(H)$。考虑相应的去等变化范畴 $\mathcal{C}'$。

% [Tutor Note: 易错预警。记号在此发生翻转：§6.1 约定 $\mathcal{C}$ 为类型 (1)、$\mathcal{C}'$ 为类型 (2)，而 §6.2 起把类型 (2) 的范畴记作 $\mathcal{C}$、把去等变化所得的类型 (1) 范畴记作 $\mathcal{C}'$。下文「对 $\mathcal{C}'$ 作等变化即恢复 $\mathrm{Rep}(H)$」正是从类型 (1) 回到类型 (2) 的方向，可据此校验。读 6.2、6.3 时务必按上下文判断 $\mathcal{C}$ 与 $\mathcal{C}'$ 各属哪一类型。]

设 $\mathcal{O}_H$ 是 $H$ 上的函数代数，即 $H$ 的正则表示。我们有同构
\[
\mathcal{O}_H = \bigoplus_{V\in\mathrm{Irrep}(H)} V\otimes V^{*},
\]
它保持两边 $H$ 的左、右作用（这里 $\mathrm{Irrep}(H)$ 是 $H$ 的不可约表示的等价类之集，对每个表示 $V$ 我们记 $V^{*}$ 为其对偶表示）。

式 (6.1) 中同构 $\alpha_V$ 的数据以及它们之间的相容性，可以简洁地概括为 $M$ 上的一个 $\mathcal{O}_H$-模结构，且与 $H$ 的左作用相容。于是 $\mathcal{C}'$ 是 $H$ 上 $H$-等变凝聚层的范畴。

该范畴等价于向量空间范畴 $\mathrm{Vect}$。事实上，我们有函子 $G : \mathrm{Vect}\to\mathcal{C}'$，把 $U\in\mathrm{Vect}$ 送到 $U\otimes\mathcal{O}_H$。其拟逆函子 $\mathcal{C}'\to\mathrm{Vect}$ 由把 $M\in\mathcal{C}'$ 送到它在 $1\in H$ 处的纤维来定义。上面描述的 $H$ 在 $\mathcal{C}'$ 上的自然作用，变为 $\mathrm{Vect}$ 上的平凡作用。相应的 $\mathrm{Vect}$ 的 $H$-等变范畴以带有 $H$ 的作用的向量空间为对象。于是，通过对 $\mathcal{C}'$ 作等变化，我们如所承诺地恢复了出发时的范畴 $\mathrm{Rep}(H)$。

% ==== 原文 PDF 第 78 页（印刷页 78）====

更具体地考虑 $H = \mathbb{Z}_2 = \{1,-1\}$ 的情形是有益的。此时 $\mathrm{Rep}(H)$ 有两个一维不可约表示：平凡表示 $I$ 与符号表示 $S$。

设 $\mathcal{C}$ 是如下类型（2）范畴：它有两个不可约对象 $A_+$ 与 $A_-$，且 $\mathrm{Rep}(\mathbb{Z}_2)$ 的作用如下：
\begin{equation}\tag{6.3}
I\star A_\pm = A_\pm,\qquad S\star A_\pm = A_\mp.
\end{equation}
$\mathcal{C}$ 的所有其他对象都是 $A_+$ 与 $A_-$ 的若干个副本的直和。

我们先假设 $A_+$ 与 $A_-$ 中有一个对应于平凡表示的优先对象，比如 $A_+$。

相应的类型（1）范畴 $\mathcal{C}'$ 是什么样子？按定义，其对象形如 $(M,(\alpha_I,\alpha_S))$，其中 $M\in\mathcal{C}$ 且
\[
\alpha_I : I\star M\simeq M,\qquad \alpha_S : S\star M\simeq M
\]
是满足下述条件的同构。首先，$I\star M$ 就是 $M$，而 $\alpha_I$ 是恒等映射 $M\to M$。其次，
\[
(\alpha_S)^2 : M = I\star M = (S\otimes S)\star M = S\star(S\star M)\to S\star M\to M
\]
是恒等映射。于是我们的数据可以记为对 $(M,\alpha)$，其中 $M\in\mathcal{C}$ 且 $\alpha = \alpha_S$ 是同构 $S\star M\to M$，满足 $\alpha^2 = \mathrm{id}$。

容易看出，若 $M$ 是 $\mathcal{C}$ 的对象，且存在满足上述性质的同构 $\alpha : S\star M\simeq M$，则 $M$ 是
\[
A := A_+\oplus A_-.
\]
的若干个副本的直和。

我们考察 $\mathcal{C}'$ 中形如 $(A,\alpha)$ 的对象。我们有 $S\star A = S(A_+)\oplus S(A_-) = A$。因此 $\alpha : A\to A$ 由两个非零标量 $\epsilon_\pm$ 决定，它们对应于 $\alpha$ 在 $A_\pm$ 上的作用（我们已假设 $A_+$ 与 $A_-$ 都不可约且不同构）。$\alpha^2 = \mathrm{id}$ 这一条件蕴含 $\epsilon_+\epsilon_- = 1$。给定两个这样的对象 $(A,(\epsilon_+,\epsilon_-))$ 与 $(A,(\epsilon'_+,\epsilon'_-))$，它们之间的同构是作用在 $A_+$ 与 $A_-$ 上的一对非零标量 $(\lambda_+,\lambda_-)$，使得 $\epsilon_+/\epsilon'_+ = \lambda_+/\lambda_-$。于是 $\lambda_+$ 由 $\lambda_-$ 决定，反之亦然。

于是范畴 $\mathcal{C}'$ 非常简单：在同构意义下只有一个不可约对象 $A$，而所有其他对象都同构于该对象的若干个副本的直和。事实上，$\mathcal{C}'$ 等价于向量空间范畴 $\mathrm{Vect}$。函子 $\mathcal{C}'\to\mathrm{Vect}$ 把 $M\in\mathcal{C}'$ 送到 $\mathrm{Hom}(A_+,M)$。

在 $SL_2$ 的几何内窥研究中，我们会遇到一个与范畴 $\mathcal{C}$ 类似的 A-膜范畴。

% [原文疑误: 同一对象在前文（p.78 前半）记作斜体 $A$、$A_\pm$，自本段起改记花体 $\mathcal{A}$、$\mathcal{A}_\pm$，而 p.79 末句「投影到 $\mathcal{A}_\pm$ 上」中的外层对象又写作斜体 $A$。这是排版/记号不一致，并非两个不同对象；译文按图像原样保留，不统一记号。]

它也有两个不可约对象 $\mathcal{A}_+$ 与 $\mathcal{A}_-$，而 $\mathrm{Rep}(\mathbb{Z}_2)$ 按公式 (6.3) 作用于它。然而，$\mathcal{A}_+$ 与 $\mathcal{A}_-$ 中没有优先对象；换言之，$\mathcal{C}$ 与 $\mathrm{Rep}(\mathbb{Z}_2)$ 之间没有典范等价。自然会问：这一模糊性在范畴 $\mathcal{C}'$ 中如何反映？

答案很清楚：范畴 $\mathcal{C}'$ 仍然等价于 $\mathrm{Vect}$，但这样的等价有两种，且无法在二者之间作出选择。相应的函子 $\mathcal{C}'\to\mathrm{Vect}$ 把 $M\in\mathcal{C}'$ 送到 $\mathrm{Hom}(\mathcal{A}_\pm,M)$。在这些等价下，$\mathcal{C}'$ 的所有对象 $(A,(\epsilon_+,\epsilon_-))$（其中 $\epsilon_+\epsilon_- = 1$）都送到一维向量空间 $\mathbb{C}$（视为 $\mathrm{Vect}$ 的对象），但同构
\[
(A,(\epsilon_+,\epsilon_-))\to(A,(\epsilon'_+,\epsilon'_-))
\]

% [Tutor Note: 物理图像。这两个互不同构的等价 $\mathcal{C}'\simeq\mathrm{Vect}$ 正是内窥情形中两个分数膜不可区分的范畴表现：$\mathrm{Hom}(\mathcal{A}_+,\cdot)$ 与 $\mathrm{Hom}(\mathcal{A}_-,\cdot)$ 分别对应奇异 Hitchin 纤维的两个分支，选取其中之一等于选取一个分支，而这恰是 §10 中 $\mathbb{Z}_2$-gerbe 所阻挠的那一步。]

% ==== 原文 PDF 第 79 页（印刷页 79）====
（如上所述由 $(\lambda_+,\lambda_-)$ 给出）映到由乘以 $\lambda_+$ 或乘以 $\lambda_-$ 给出的同构 $\mathbb{C}\to\mathbb{C}$。

于是，两个对象 $\mathcal{A}_+$ 与 $\mathcal{A}_-$ 各自对应一个特定的等价 $\mathcal{C}'\simeq\mathrm{Vect}$，但既然我们无法在 $\mathcal{A}_+$ 与 $\mathcal{A}_-$ 之间作出选择，我们也就无法在这两个等价之间作出取舍（容易看出，作为函子 $\mathcal{C}'\to\mathrm{Vect}$，这些等价彼此不同构）。当我们过渡到类型（1）的范畴 $\mathcal{C}'$ 时，正是这一点取代了类型（2）范畴 $\mathcal{C}$ 中 $\mathcal{A}_+$ 与 $\mathcal{A}_-$ 之间的模糊性。

范畴等价 $\mathcal{C}'\simeq\mathrm{Vect}$ 有两种不同取法，这确实看起来比「范畴 $\mathcal{C}$ 有两个不可区分的对象」这一更具体的说法更微妙、更玄奥。这很好地说明了为什么从实用角度看，处理类型（2）的范畴往往比处理类型（1）的范畴更好。

然而要点在于，这两种描述彼此等价。为使自己确信这一点，看清如何从 $\mathcal{C}'$ 恢复 $\mathcal{C}$ 是有启发的。我们对 $\mathcal{C}'$ 施行等变化程序。这样我们定义一个新范畴 $\mathcal{C}''$，其对象是对 $(M,\mu)$，其中 $M$ 是 $\mathcal{C}'$ 的对象，而 $\mu = \mu_{-1}$ 是 $M$ 与如下新对象 $M'$ 之间的同构：$M'$ 由把对应于 $-1\in\mathbb{Z}_2$ 的函子作用于 $M$ 得到。注意 $\mu_1 = \mathrm{id} : M\to M$，因此为满足 $\mathbb{Z}_2$ 的关系，我们必须有
\[
\mu^2 = \mathrm{id}.
\]
若 $M = (\mathcal{A},(\epsilon_+,\epsilon_-))$，则 $M' = (\mathcal{A},(-\epsilon_+,-\epsilon_-))$，而 $M$ 与 $M'$ 之间恰有两个满足 $\mu^2 = \mathrm{id}$ 的同构 $\mu_\pm$：一个在 $\mathcal{A}_\pm\subset\mathcal{A}$ 上作用为 $\pm1$，另一个作用为 $\mp1$。于是我们看到，$\mathcal{C}''$ 中有两个不同构的对象，它们在典范等价 $\mathcal{C}''\simeq\mathcal{C}$ 下分别对应于 $\mathcal{A}_+$ 与 $\mathcal{A}_-$。事实上，把 $\frac{1}{2}(1+\mu_\pm)$ 设想为 $\mathcal{A}$ 中到 $\mathcal{A}_\pm$ 上的投影算子是有启发的。

% [Tutor Note: 推导补足。这里用到了幂等态射的分裂：$\mu_\pm^2=\mathrm{id}$ 意味着 $\frac{1}{2}(1+\mu_\pm)$ 是幂等态射，其像给出 $\mathcal{A}$ 的两个直和项；在特征 $0$ 的线性范畴中幂等态射自动分裂，故 $\mathcal{C}''$ 的两个不同构对象恰是这两个像。这也解释了为何从 $\mathcal{C}'$（唯一不可约对象 $\mathcal{A}$）作等变化能还原 $\mathcal{C}$ 的两个不可约对象。]

\subsection{Hecke 本征层}

我们之所以如此详细地讨论这些微妙之处，是因为与经典自守形式理论中通常的 Hecke 本征函数概念最接近的范畴是类型（1）的范畴；因此要完全理解与经典自守形式理论的联系，我们必须经过类型（1）的范畴。然而，第 5 节所考虑的 A-膜范畴自然地是类型（2）的范畴。所以我们必须在两类范畴之间建立联系。

我们先回顾几何 Langlands 纲领中所用的 Hecke 本征层的传统定义（例如见 [BD] 或 [F1] 第 6.1 节）。它们是 $\Bun_G$ 上的 D-模，即曲线 $C$ 上 $G$-丛的模叠上的 D-模，满足 Hecke 本征对象性质。为更精确地解释这一点，回忆对偶群 ${}^LG$ 的每个有限维表示 $V$ 都给出一个 Hecke 函子 $H_V$，作用从 $\Bun_G$ 上的 D-模范畴到 $C\times\Bun_G$ 上的 D-模范畴。这里我们不回顾这些函子的定义，请读者参阅 [BD] 与 [F1]。我们指出，如 [KW] 所解释的，这些函子与第 5.3 节讨论的 't Hooft 算子密切相关。

% ==== 原文 PDF 第 80 页（印刷页 80）====

设 $\mathcal{E}$ 是 $C$ 上的平坦 ${}^LG$-丛。\footnote{这里为简单起见我们考虑非分歧情形，但该定义容易推广到分歧情形；只需去掉 $\mathcal{E}$ 的分歧点。}以 $\mathcal{E}$ 为「本征值」（eigenvalue）的 Hecke 本征层，按定义是一族数据
\begin{equation}\tag{6.4}
(\mathcal{F},(\alpha_V)_{V\in\mathrm{Rep}({}^LG)}),
\end{equation}
其中 $\mathcal{F}$ 是 $\Bun_G$ 上的 D-模，而 $(\alpha_V)$ 是一族同构
\begin{equation}\tag{6.5}
\alpha_V : H_V(\mathcal{F})\xrightarrow{\sim} V_{\mathcal{E}}\boxtimes\mathcal{F},
\end{equation}
其中
\[
V_{\mathcal{E}} = \mathcal{E}\underset{{}^LG}{\times} V
\]
是与 $V$ 相联系的 $C$ 上的平坦向量丛。

这些同构必须满足自然的相容性条件：关于 (6.5) 左端 Hecke 函子 $H_V$ 的复合，以及关于 (6.5) 右端表示的张量积，此外还有自然的结合性条件。\footnote{我们将在第 8.2 节看到，还需对同构 $\alpha_V$ 施加一个额外的交换性条件。}

设 $Aut_{\mathcal{E}}$ 是所有本征值为 $\mathcal{E}$ 的 Hecke 本征层构成的范畴。为与上一节所研究的类型（1）范畴建立联系，固定一点 $x\in C$。则 Hecke 算子 $H_V$ 到 $x$ 的限制是从 $\Bun_G$ 上 D-模范畴到自身的函子 $H_{V,x}$。对一个 Hecke 本征层 (6.4)，把同构 $\alpha_V$ 限制到 $x$，我们得到一族相容的同构
\begin{equation}\tag{6.6}
\alpha_{V,x} : H_{V,x}(\mathcal{F})\xrightarrow{\sim} V_{\mathcal{E},x}\otimes\mathcal{F}.
\end{equation}
这里
\[
V_{\mathcal{E},x} = \mathcal{E}_x\underset{{}^LG}{\times} V,
\]
其中 $\mathcal{E}_x$ 是 $\mathcal{E}$ 在 $x$ 处的纤维，是一个同构于 $V$ 的向量空间。$\mathcal{F}$ 与 $(\alpha_{V,x})_{V\in\mathrm{Rep}({}^LG)})$ 的数据，正是我们上面在去等变化范畴的定义中所用的那类数据。

% [原文疑误: 此处 $(\alpha_{V,x})_{V\in\mathrm{Rep}({}^LG)})$ 多了一个右括号；与 (6.4)、(6.9) 的写法对照，应为 $(\alpha_{V,x})_{V\in\mathrm{Rep}({}^LG)}$。译文照原样保留该括号，以存原文之旧。]

因此 Hecke 本征层 (6.4) 是这种类型的对象，只不过我们拥有的不是一族 Hecke 同构 (6.1)，而是由曲线 $C$ 上的点参数化的一整族这样的同构族。条件 (6.5) 实际上比所有 $x\in C$ 的同构 (6.6) 的数据包含多得多的信息，因为公式 (6.5) 还描述了 Hecke 算子的「本征值」$V_{\mathcal{E},x}$ 对 $x$ 的依赖：它们按局部系统 $\mathcal{E}$ 「变化」。

若我们只施加条件 (6.6)（对某个固定的 $x\in C$），则相应的范畴带有群 ${}^LG$ 的作用，或者更确切地说，带有它经 ${}^LG$-torsor $\mathcal{E}_x$ 的扭转的作用，即 ${}^LG_{\mathcal{E}_x} = \mathcal{E}_x\underset{{}^LG}{\times}{}^LG$ 的作用。该作用按公式 (6.2) 定义。另一方面，若我们施加完整的 Hecke 条件 (6.1)，则群 ${}^LG_{\mathcal{E}_x}$ 不再作用。此时我们有的是 $\mathcal{E}$ 的（整体）自同构群的作用，即 $\Gamma(C,{}^LG_{\mathcal{E}})$ 的作用，其中 ${}^LG_{\mathcal{E}} = \mathcal{E}\underset{{}^LG}{\times}{}^LG$。

% [Tutor Note: 易错预警。须区分局部与整体两种对称性：只固定一点 $x$ 时，扭转后的群 ${}^LG_{\mathcal{E}_x}$ 仍然作用，但它依赖于 $\mathcal{E}_x$ 的平凡化；施加完整条件 (6.5)（对曲线 $C$ 上所有 $x$ 同时成立）后，残余对称性只剩常数自同构群 $\Gamma(C,{}^LG_{\mathcal{E}})$。换言之，(6.6) 对单个 $x$ 的数据远少于 (6.5)：后者还规定了本征值如何随 $x$ 按局部系统 $\mathcal{E}$ 变化。]

假设我们的局部系统 $\mathcal{E}$ 的自同构群是平凡的。特别地，这意味着 ${}^LG$ 的中心平凡，从而 ${}^LG$ 是伴随型半单群。在这种情况下，预期存在唯一的不可约 D-模
% ==== 原文 PDF 第 81 页（印刷页 81）====
$\mathcal{F}$ 满足 Hecke 性质 (6.1)，而所有其他满足该性质的 D-模都是它的若干个副本的直和。若是如此，则对不可约 D-模 $\mathcal{F}$，构造 (6.5) 中的同构 $\alpha_V$ 本质上只有一种方式；这就是为什么在此情形下人们通常省略 $(\alpha_V)$ 的数据。

然而，当 $\mathcal{E}$ 有非平凡的自同构群时，这些数据就变得重要了：
\[
\Gamma := \Gamma(C,{}^LG_{\mathcal{E}}).
\]
于是我们像公式 (6.2) 那样定义群 $\Gamma$ 在范畴 $Aut_{\mathcal{E}}$ 上的作用。也就是说，给定范畴 $Aut_{\mathcal{E}}$ 中如公式 (6.4) 所示的对象 $(\mathcal{F},(\alpha_V))$，我们构造一个新对象 $(\mathcal{F},(\alpha_V^s))$，其中
\[
\alpha_V^s = (s\otimes1)\circ\alpha_V.
\]
这个新对象可能不同构于旧对象（即使同构，也可能以不同的方式同构于它）。

\subsection{内窥例中的 Hecke 本征层范畴}

下面讨论内窥例子中的 Hecke 本征层范畴，此时 $G = SL_2$、${}^LG = SO_3$、$\Gamma = \mathbb{Z}_2$。我们预期在此情形下，任何满足 Hecke 本征层性质的 D-模都是某个 D-模的若干个副本的直和，该 D-模我们记作 $\mathcal{F}$。如第 2.3 节与第 4 节所解释的，D-模 $\mathcal{F}$ 对应于 $\mathcal{M}_H(G)$ 的某个奇异纤维上的 A-膜 $\mathcal{A}$，它是关于 't Hooft 算子的磁本征膜。如第 5 节所解释的，这个 A-膜是可约的：
\[
\mathcal{A} = \mathcal{A}_+\oplus\mathcal{A}_-,
\]
其中 A-膜 $\mathcal{A}_\pm$ 不可约。此外，它们之中没有一个更优先（关于这一点的更详细讨论见第 10 节）。实际上，正是由于这种模糊性，我们先前对这些 A-膜使用记号 $\mathcal{A}_1$ 与 $\mathcal{A}_2$，以强调它们并不典范地对应于 B-膜 $\mathcal{B}_+$ 与 $\mathcal{B}_-$。但从现在起我们将使用记号 $\mathcal{A}_\pm$。

因此我们预期 D-模 $\mathcal{F}$ 也是可约的：
\[
\mathcal{F} = \mathcal{F}_+\oplus\mathcal{F}_-,
\]
且每个 $\mathcal{F}_\pm$ 不可约。我们还预期二者之中没有哪一个比另一个更优先。

回忆上面定义的本征层概念包含对 $SO_3$ 的所有表示 $V$ 的同构 $\alpha_V$。利用与张量积结构的相容性，我们发现一切都由 $SO_3$ 的伴随表示决定，我们如前记之为 $W$。因此 Hecke 本征层可以看作一对 $(\mathcal{F},\alpha)$，其中
\begin{equation}\tag{6.7}
\alpha : H_W(\mathcal{F})\xrightarrow{\sim} W_{\mathcal{E}}\boxtimes\mathcal{F}.
\end{equation}
在内窥情形，我们的 $SO_3$-局部系统 $\mathcal{E}$ 的结构群约化到子群
\[
O_2 = \mathbb{Z}_2\ltimes\mathbb{C}^\times
\]
% ==== 原文 PDF 第 82 页（印刷页 82）====
（但不约化到更小的子群）。记 $U$ 为 $O_2$ 的定义二维表示。则 $\det U$ 是由同态 $O_2\to\mathbb{Z}_2$ 诱导的一维符号表示。作为 $O_2\times\mathbb{Z}_2$ 的表示，我们有
\[
W = (\det U\otimes I)\oplus(U\otimes S),
\]
其中 $\mathbb{Z}_2$ 是 $O_2$ 在 $SO_3$ 中的中心化子（它等于 $O_2$ 的中心），且如前所述，$S$ 是 $\mathbb{Z}_2$ 的符号表示，$I$ 是 $\mathbb{Z}_2$ 的平凡表示。因此我们得到相应局部系统的如下分解：
\begin{equation}\tag{6.8}
W_{\mathcal{E}} = (\det U_{\mathcal{E}}\otimes I)\oplus(U_{\mathcal{E}}\otimes S).
\end{equation}

% [Tutor Note: 物理图像。$O_2=\mathbb{Z}_2\ltimes\mathbb{C}^\times$ 有两个连通分支，$\det U$ 正是探测该分支的符号特征，因此 $\det U_{\mathcal{E}}$ 是 2 阶线丛给出的局部系统；$SO_3$ 的伴随表示 $W$ 限制到 $O_2\times\mathbb{Z}_2$ 后分裂为 $\det U\otimes I$ 与 $U\otimes S$ 两项，这两项就是 (6.11) 中两个本征值的来源，并与 §5 公式 (5.32) 中 't Hooft 算子对分数膜的置换作用一致。]

我们可以用非平凡元素 $\tau\in\mathbb{Z}_2 = \Gamma$ 的作用来扭转同构 $\alpha$，该元素属于我们局部系统的自同构群。这样我们得到一个新的 Hecke 本征层，即一对 $(\mathcal{F},\alpha')$，其中 $\alpha' = (\tau\otimes1)\circ\alpha$。$Aut_{\mathcal{E}}$ 的对象 $(\mathcal{F},\alpha)$ 与 $(\mathcal{F},\alpha')$ 同构，但并非典范同构。事实上存在两个自然的同构，分别等于 $\mathcal{F}_\pm$ 上的 $\pm1$ 或 $\mp1$，且没有自然的方式在二者之间作出选择。

我们预期 $Aut_{\mathcal{E}}$ 中任何 Hecke 本征层对象都同构于 $(\mathcal{F},\alpha)$ 的若干个副本的直和。如第 6.2 节讨论的玩具模型那样，这意味着范畴 $Aut_{\mathcal{E}}$ 等价于向量空间范畴，但有两种不同方式，对应于 $\mathcal{F}_+$ 与 $\mathcal{F}_-$ 的选取。

\subsection{分数 Hecke 本征层}

接下来，我们引入类型（2）的 Hecke 本征层范畴。

再次假设给定曲线 $C$ 上的 ${}^LG$-局部系统 $\mathcal{E}$，设 $\Gamma$ 是它的自同构群。为简化下面的讨论，我们通过选取一点 $x\in C$ 并选取 $\mathcal{E}$ 在 $x$ 处纤维 $\mathcal{E}_x$ 的一个平凡化，把 $\Gamma$ 等同于 ${}^LG$ 的一个子群。（这使我们能给 $\mathcal{E}$ 赋予一个同态 $\pi_1(C,x)\to{}^LG$。群 $\Gamma$ 于是可以定义为其像的中心化子。）

假设给定 $\Bun_G$ 上 D-模范畴的一个阿贝尔子范畴 $\mathcal{C}$，它带有张量范畴 $\mathrm{Rep}(\Gamma)$ 的作用。换言之，对每个 $R\in\mathrm{Rep}(\Gamma)$ 我们有一个函子
\[
M\mapsto R\star M,
\]
且这些函子的复合方式与 $\mathrm{Rep}(\Gamma)$ 上的张量积结构相容。「本征值」为 $\mathcal{E}$ 的类型（2）Hecke 本征层范畴将以如下数据为对象：
\begin{equation}\tag{6.9}
(\mathcal{F},(\alpha_V)_{V\in\mathrm{Rep}({}^LG)}),
\end{equation}
其中 $\mathcal{F}$ 是 $\mathcal{C}$ 的对象，而 $\alpha_V$ 是下面定义的同构。

记 $\mathrm{Res}_\Gamma(V)$ 为 ${}^LG$ 的表示 $V$ 到 $\Gamma$ 的限制。若 $\mathrm{Rep}(\Gamma)$ 是半单范畴（例如当 $\Gamma$ 是有限群时就是如此），则我们得到分解
\[
\mathrm{Res}_\Gamma(V) = \bigoplus_i F_i\otimes R_i,
\]
% ==== 原文 PDF 第 83 页（印刷页 83）====
其中 $R_i$ 是 $\Gamma$ 的不可约表示，$F_i$ 是 ${}^LG$ 中 $\Gamma$ 的中心化子的相应表示。用 $\mathcal{E}$ 扭转，我们得到 $C$ 上带有 $\Gamma$ 的交换作用的一个局部系统 $(\mathrm{Res}_\Gamma(V))_{\mathcal{E}}$，它分解如下：
\[
(\mathrm{Res}_\Gamma(V))_{\mathcal{E}} = \bigoplus_i (F_i)_{\mathcal{E}}\otimes R_i.
\]
注意，由于 $\Gamma$ 是 $\mathcal{E}$ 的自同构群，$\mathcal{E}$ 的结构群约化到 $\Gamma$ 在 $G$ 中的中心化子，而 $F_i$ 是该中心化子的表示。因此 $F_i$ 可以用 $\mathcal{E}$ 扭转，所得的 $C$ 上的局部系统（或平坦向量丛）记作 $(F_i)_{\mathcal{E}}$。

% [原文疑误: $\mathcal{E}$ 是 ${}^LG$-局部系统，其结构群约化后应是 $\Gamma$ 在 ${}^LG$（而非 $G$）中的中心化子；本页上一行正写作「$\Gamma$ 在 ${}^LG$ 中的中心化子的相应表示」。此处「in $G$」当为「in ${}^LG$」之误。译文照原样保留 $G$。]

同构 $\alpha_V$ 形如
\begin{equation}\tag{6.10}
\alpha_V : H_V(M)\xrightarrow{\sim}(\mathrm{Res}_\Gamma(V))_{\mathcal{E}}\star M = \bigoplus_i(F_i)_{\mathcal{E}}\boxtimes(R_i\star M),\qquad M\in\mathcal{C},
\end{equation}
且它们必须在显然的意义下相容。我们把对象为 (6.9) 且满足上述条件的范畴记作 $Aut'_{\mathcal{E}}$。\footnote{如我们将在第 8.2 节看到的，还需要对同构 $\alpha_V$ 施加一个额外的交换性条件。}

% [Tutor Note: 易错预警。比较 (6.5) 与 (6.10) 须逐项对应且注意两端的不同：(6.5) 左端是 $H_V(\mathcal{F})$、右端是 $V_{\mathcal{E}}\boxtimes\mathcal{F}$；类型 (2) 的 (6.10) 左端却是 $H_V(M)$（$M\in\mathcal{C}$），右端把本征值局部系统吸收进 $(\mathrm{Res}_\Gamma(V))_{\mathcal{E}}$，再用 $\star$ 保留 $\mathrm{Rep}(\Gamma)$ 的幺半作用。$V_{\mathcal{E}}\boxtimes\mathcal{F}$ 对应的是 $\bigoplus_i(F_i)_{\mathcal{E}}\boxtimes(R_i\star M)$，而不是 $(\mathrm{Res}_\Gamma(V))_{\mathcal{E}}\boxtimes M$。]

类型（2）的 Hecke 本征层范畴的优点是比类型（1）的范畴更具体，并且它更贴近我们在第 5 节所发现的 A-膜与 B-膜范畴的结构。

然而，范畴 $Aut_{\mathcal{E}}$ 可以沿第 6.1 节的思路，通过去等变化程序从 $Aut'_{\mathcal{E}}$ 重构出来。反之，对 $Aut_{\mathcal{E}}$ 施行等变化程序，我们得到一个等价于 $Aut'_{\mathcal{E}}$ 的范畴。

在我们的几何内窥主例子中，范畴 $Aut'_{\mathcal{E}}$ 是什么样子？在此情形，范畴 $\mathcal{C}$ 应当有两个不可约对象 $\mathcal{F}_+$ 与 $\mathcal{F}_-$，它们是 $\Bun_G$ 上对应于分数 A-膜 $\mathcal{A}_+$ 与 $\mathcal{A}_-$ 的 D-模。范畴 $\mathrm{Rep}(\mathbb{Z}_2)$ 对它们的作用如下：$\Gamma = \mathbb{Z}_2$ 的符号表示 $S$ 置换它们，
\[
S\star\mathcal{F}_\pm = \mathcal{F}_\mp,
\]
而平凡表示 $I$ 恒等地作用。

由于 $SO_3$ 的表示范畴由伴随表示 $W$ 生成，只需对 $SO_3$ 的伴随表示 $W$ 表述 Hecke 性质 (6.10)。它读作
\begin{equation}\tag{6.11}
\begin{aligned}
H_W(\mathcal{F}_+) &\simeq (\det U_{\mathcal{E}}\boxtimes\mathcal{F}_+)\oplus(U_{\mathcal{E}}\boxtimes\mathcal{F}_-),\\
H_W(\mathcal{F}_-) &\simeq (\det U_{\mathcal{E}}\boxtimes\mathcal{F}_-)\oplus(U_{\mathcal{E}}\boxtimes\mathcal{F}_+),
\end{aligned}
\end{equation}
其中 $\det U_{\mathcal{E}}$ 与 $U_{\mathcal{E}}$ 是公式 (6.8) 中定义的 $W_{\mathcal{E}}$ 的直和项。这与公式 (5.32) 给出的 't Hooft 算子对 A-膜的作用相符。由于该公式描述的是分数膜 $\mathcal{A}_\pm$ 的行为，我们把由公式 (6.10) 与 (6.11) 表达的性质称为分数 Hecke 性质，把相应的 D-模称为分数 Hecke 本征层。另一方面，我们把通常的 Hecke 性质 (6.5) 称为正则 Hecke 性质，把满足它的 D-模称为正则 Hecke 本征层（这一术语的理由是：这样的本征层对应于群 $\Gamma$ 的正则表示）。

% ==== 原文 PDF 第 84 页（印刷页 84）====

\subsection{其他例子}

本节我们考察分数 Hecke 本征层范畴的其他例子。首先考虑阿贝尔情形，即 $G = \mathbb{C}^\times$ 或任意环面。然后考虑 ${}^LG$-局部系统的情形（其中 $G$ 是任意约化群），其自同构群为 ${}^LG$ 的中心。最后我们讨论几何 Eisenstein 级数。

\subsubsection{\texorpdfstring{$\mathbb{C}^\times$}{C×} 的情形}

设 $G = \mathbb{C}^\times$，于是 ${}^LG = \mathbb{C}^\times$。设 $\mathcal{E}$ 是 $C$ 上的 $\mathbb{C}^\times$-局部系统，我们把它视为秩一局部系统。范畴 $Aut_{\mathcal{E}}$ 的对象是数据 $(\mathcal{F},\alpha)$，其中 $\mathcal{F}$ 是 $\mathrm{Pic}$（$C$ 的 Picard 概形）上的 D-模，且
\begin{equation}\tag{6.12}
\alpha : H_1(\mathcal{F})\xrightarrow{\sim}\mathcal{E}\boxtimes\mathcal{F}.
\end{equation}
这里 $H_1$ 是对应于恒等特征 $\mathbb{C}^\times\to\mathbb{C}^\times$ 的 Hecke 函子，
\[
H_1(\mathcal{F}) = p^{*}(\mathcal{F}),
\]
其中
\[
p : C\times\mathrm{Pic}\to\mathrm{Pic},\qquad (x,\mathcal{L})\mapsto\mathcal{L}(-x).
\]
这一 Hecke 性质自动蕴含对应于 $\mathbb{C}^\times$ 的其他特征的 Hecke 函子 $H_n$（$n\in\mathbb{Z}$）的 Hecke 性质。由 [La4, Ro] 给出的 D-模 Fourier--Mukai 等价蕴含 $\mathcal{F}$ 是 $\mathrm{Pic}$ 上某个特定平坦线丛的若干个副本的直和。今后我们用 $\mathcal{F}$ 记这个平坦线丛。

% [Tutor Note: 推导补足。D-模的 Fourier--Mukai 等价把 Hecke 函子 $H_1$ 与「张量积以 Poincaré 线丛」互换；于是本征条件 $\alpha:H_1(\mathcal{F})\simeq\mathcal{E}\boxtimes\mathcal{F}$ 化为对 Fourier--Mukai 变换后的对象的秩一条件，再作逆变换即得 $\mathcal{F}$ 是某个平坦线丛的若干个副本的直和。原文用「Fourier--Mukai 等价」一句带过的正是这一步。]

由于 $\mathrm{Pic}$ 是其各分支 $\mathrm{Pic}_n$（$n\in\mathbb{Z}$，对应于次数 $n$ 的线丛）的不交并，$\mathrm{Pic}$ 上的每个 D-模 $\mathcal{F}$ 都分解为直和
\begin{equation}\tag{6.13}
\mathcal{F} = \bigoplus_{n\in\mathbb{Z}}\mathcal{F}_n.
\end{equation}

任何秩一局部系统 $\mathcal{E}$ 的自同构群都是 $\mathbb{C}^\times$。因此我们有 $\mathbb{C}^\times$ 在范畴 $Aut_{\mathcal{E}}$ 上的作用。它由公式
\[
\lambda\cdot(\mathcal{F},\alpha)\mapsto(\mathcal{F},\lambda\alpha),\qquad \lambda\in\mathbb{C}^\times
\]
给出。

现在描述相应的分数 Hecke 本征层范畴。设 $\mathcal{C}$ 是这样的范畴：其对象是 Hecke 本征层 $\mathcal{F}$ 的分解 (6.13) 中出现的 D-模 $\mathcal{F}_n$（$n\in\mathbb{Z}$）的直和。这些是分支 $\mathrm{Pic}_n\subset\mathrm{Pic}$ 上的平坦向量丛。范畴 $\mathrm{Rep}(\mathbb{C}^\times)$ 在 $\mathcal{C}$ 上的作用由公式
\[
[m]\star\mathcal{F}_n = \mathcal{F}_{n+m}
\]
给出，其中 $[m]$ 表示对应于特征 $\mathbb{C}^\times\to\mathbb{C}^\times$（$a\mapsto a^m$）的一维表示。

分数 Hecke 本征层范畴 $Aut'_{\mathcal{E}}$ 的对象于是由数据 $(\mathcal{K},\alpha)$ 给出，其中 $\mathcal{K}\in\mathcal{C}$ 且
\begin{equation}\tag{6.14}
\alpha : H_1(\mathcal{K})\xrightarrow{\sim}\mathcal{E}\boxtimes([1]\star\mathcal{K}).
\end{equation}
同构 $\alpha$ 给出唯一的同构
\[
\alpha_n : H_n(\mathcal{K})\xrightarrow{\sim}\mathcal{E}^{\otimes n}\boxtimes([n]\star\mathcal{K}),\qquad n\in\mathbb{Z},
\]
满足所需性质。

% ==== 原文 PDF 第 85 页（印刷页 85）====

注意，由公式 (6.12) 给出的 $\mathcal{F}$ 的正则 Hecke 性质等价于各 $\mathcal{F}_n$（$n\in\mathbb{Z}$）的分数 Hecke 性质。因此对 $\mathcal{K} = \mathcal{F}_m$ 我们得到同构 (6.14)，从而对 $\mathcal{C}$ 的所有对象都成立。

这个分数 Hecke 本征层范畴的例子应当与我们来自内窥的 $G = SL_2$ 主例子相对照。在后一情形，自同构群是有限群 $\mathbb{Z}_2$。因此正则 Hecke 本征层的分解是有限的，$\mathcal{F} = \mathcal{F}_+\oplus\mathcal{F}_-$，其直和项由 $\mathbb{Z}_2$ 的不可约表示标记，尽管这种标记并非典范。在另一例子中，自同构群是 Lie 群 $\mathbb{C}^\times$。因此正则 Hecke 本征层分解为无穷多个直和项之和（6.13），这些直和项这次是由 $\mathbb{C}^\times$ 的不可约表示典范地标记的。

一个重要的差别在于：在内窥例子中，分数 Hecke 本征层范畴 $\mathcal{C}$ 与范畴 $\mathrm{Rep}(\mathbb{Z}_2)$ 之间没有典范等价（换言之，我们无法区分 $\mathcal{F}_+$ 与 $\mathcal{F}_-$）；而在另一例子中，我们把它典范地等同于 $\mathrm{Rep}(\mathbb{C}^\times)$（确实，有一个支撑在次数 0 的线丛构成的分支 $\mathrm{Pic}_0$ 上的典范对象 $\mathcal{F}_0$）。其原因将在第 10 节讨论。

\subsubsection{任意环面}

这个分数 Hecke 本征层范畴的例子直接推广到 $G$ 为任意环面的情形。此时相应的模空间 $\Bun_T$ 分解为分支 $\Bun_{T,\chi}$ 的不交并，其中 $\chi$ 取遍 ${}^LT$ 的特征格 $\check{P}$。因此正则 Hecke 本征层 $\mathcal{F}$ 分解为直和
\begin{equation}\tag{6.15}
\mathcal{F} = \bigoplus_{\chi\in\check{P}}\mathcal{F}_\chi,
\end{equation}
其中 $\mathcal{F}_\chi$ 支撑在 $\Bun_{T,\chi}$ 上。D-模 $\mathcal{F}_\chi$ 于是构成分数 Hecke 本征层范畴的构件。它们满足分数 Hecke 性质
\begin{equation}\tag{6.16}
H_\lambda(\mathcal{F}_\chi)\simeq\mathcal{F}_{\lambda+\chi} = [\lambda]\star\mathcal{F}_\chi,
\end{equation}
其中 $H_\lambda$ 是对应于特征 $\lambda\in\check{P}$ 的 Hecke 函子。范畴 $Aut'_{\mathcal{E}}$ 的对象是族 $(\mathcal{K},(\alpha_\lambda)_{\lambda\in\check{P}})$，其中 $\mathcal{K}$ 是 $\mathcal{F}_\chi$ 的若干个副本的直和（带某些重数），而诸 $\alpha_\lambda$ 构成一族相容的同构 (6.16)。

\subsubsection{中心作为自同构群}

这对任意约化群 $G$ 有类似的陈述。即，设 ${}^LZ$ 是对偶群 ${}^LG$ 的中心。假设 $\mathcal{E}$ 是 $C$ 上的 ${}^LG$-局部系统，其自同构群恰为 ${}^LZ$。一般的局部系统具有这一性质。在此情形我们预期存在唯一的正则 Hecke 本征层 $\mathcal{F}$，其本征值为 $\mathcal{E}$，且在 $\Bun_G$ 的每个连通分支上不可约（而任何其他 Hecke 本征层都是 $\mathcal{F}$ 的若干个副本的直和）。$\Bun_G$ 的连通分支恰由 ${}^LZ$ 的特征格 $\check{P}$ 标记，所以我们处于与上面相同的情形：层 $\mathcal{F}$ 分解为
\[
\mathcal{F} = \bigoplus_{\chi\in\check{P}}\mathcal{F}_\chi,
\]
% ==== 原文 PDF 第 86 页（印刷页 86）====
其中 $\mathcal{F}_\chi$ 支撑在对应于 $\chi$ 的分支上。D-模 $\mathcal{F}_\chi$ 于是构成分数 Hecke 本征层范畴的构件，并满足性质 (6.16)。范畴 $Aut'_{\mathcal{E}}$ 的对象则如 $G$ 为环面的情形那样定义。我们预期该范畴等价于范畴 $\mathrm{Rep}({}^LZ)$。这与对应于 $\mathcal{E}$ 的 ${}^LG$-局部系统模叠上该点处的 B-膜范畴的结构相符，后者也等价于 $\mathrm{Rep}({}^LZ)$，因为假设 ${}^LZ$ 是 $\mathcal{E}$ 的自同构群（见第 2.2 节）。

\subsubsection{几何 Eisenstein 级数}

作为最后一个例子，我们借助 [La3, BG] 中几何 Eisenstein 级数的构造，描述非阿贝尔群的分数 Hecke 本征层范畴。更确切地说，设 $G$ 是约化代数群，$T$ 是其极大环面。则 ${}^LT$ 是 ${}^LG$ 中的极大环面。设 $\mathcal{E}$ 是 ${}^LT$-局部系统，我们把它视为 ${}^LG$-局部系统。假设 $\mathcal{E}$ 是一般的，即对应于 ${}^LG$ 的根的各秩一局部系统都非平凡。则 $\mathcal{E}$ 的自同构群等于 ${}^LT$。在此情形，存在 $\Bun_G$ 上本征值为 $\mathcal{E}$ 的 Hecke 本征层 $\mathcal{F}_G$ 的几何构造 [La3, BG]，它从关于 $\mathcal{E}$（视为 ${}^LT$-局部系统）的 Hecke 本征层 $\mathcal{F}_T$ 出发。

按构造，分解 (6.15) 给出 $\mathcal{F}_G$ 的一个分解
\begin{equation}\tag{6.17}
\mathcal{F}_G = \bigoplus_{\chi\in\check{P}}\mathcal{F}_{G,\chi}.
\end{equation}

D-模 $\mathcal{F}_{G,\chi}$ 于是构成对应于 $\mathcal{E}$ 的分数 Hecke 本征层范畴的构件。我们定义范畴 $\mathcal{C}$ 为这样的范畴：其对象是 D-模 $\mathcal{F}_{G,\chi}$（$\chi\in\check{P}$）的直和。范畴 $\mathrm{Rep}({}^LT)$ 按公式
\[
[\lambda]\star\mathcal{F}_{G,\chi} = \mathcal{F}_{G,\chi+\lambda}
\]
作用于它。

相应的范畴 $Aut'_{\mathcal{E}}$ 的对象是族 $(\mathcal{K},(\alpha_\lambda)_{\lambda\in\check{P}})$，其中 $\mathcal{K}\in\mathcal{C}$，而诸 $\alpha_\lambda$ 构成一族相容的同构（见公式 (6.10)）
\[
\alpha_\lambda : H_V(\mathcal{F}_{G,\chi})\xrightarrow{\sim}\bigoplus_\mu V(\mu)\otimes([\lambda]\star\mathcal{F}_{G,\chi}) = \bigoplus_\mu V(\mu)\otimes\mathcal{F}_{G,\chi+\mu}.
\]

% [存疑: 此式左端为 $H_V$、下标却为 $\lambda$，而 (6.10) 中该同构的下标应与表示 $V$ 一致；同时右端求和项用 $[\lambda]\star\mathcal{F}_{G,\chi}$，按 (6.16) 的类比或应为 $[\mu]$。原文如此，无法据图像判定何者为笔误，故照原样保留。]

这里 $V$ 是 ${}^LG$ 的表示，$V(\mu)$ 是 $V$ 的权 $\mu$ 子空间，于是我们有
\[
\mathrm{Res}_{{}^LT}(V) = \bigoplus_\mu V(\mu).
\]

注意，在 $G = T$ 的阿贝尔情形，分解 (6.15) 有几何来源：它对应于 $\Bun_T$ 分裂为连通分支之并。然而对非阿贝尔 $G$，相应 Eisenstein 层的分解 (6.17) 并不以任何显然的方式直接联系到底层 $G$-丛模叠的几何。这与内窥情形中所发生的情况类似。

与内窥例子类比，利用 Hitchin 纤维化的镜像对称显式构造对应于 Eisenstein 层 $\mathcal{F}_G$ 与 $\mathcal{F}_{G,\chi}$ 的 A-膜将是有意义的。在 B-模型一侧，这对应于一个更复杂的奇点，其自同构群是连续群（${}^LG$ 的极大环面），而不是内窥例子中那样的有限群（这特别
% ==== 原文 PDF 第 87 页（印刷页 87）====
反映在 Hitchin 纤维现在有无穷多个不可约分支这一事实中）。因此，在此情形镜像对称的分析变得更加微妙。我们计划在别处更详细地讨论这一点。

% [跨页接续: 原文第 87 页。§6 至此结束；自本页「7. The Classical Story」标题起属下一分片 s07-classical]
